% Introduction. Numbers are quoted from the results sections; sources are given there.
Bettors overpay for unlikely outcomes. This favourite--longshot bias was first documented at the
racetrack \citep{griffith1949}. It has since been found in most betting markets, and it is one of
the most robust anomalies in the economics of risk \citep{thalerziemba1988,sauer1998,ottaviani2008}.
Event-contract exchanges now give the old question a new setting. In the United States, Kalshi
settles hundreds of thousands of contracts a month on sports, weather, politics, entertainment and
what speakers will say. A large retail brokerage routes its customers' orders to it
\citep{robinhood2025}. Polymarket runs a comparable venue outside the United States and records
every trade on a public blockchain. These markets are often read as probability forecasts
\citep{wolferszitzewitz2004}, so whether their prices are biased, and for whom, matters beyond the
traders themselves.

We ask three questions.
\begin{enumerate}
  \item How large is the favourite--longshot bias on a regulated event-contract exchange, and how
  does it vary across event categories, horizons, contract structures, liquidity and time?
  \item Who buys the longshots, and who bears the losses: takers or makers, small or large
  orders, new or experienced accounts?
  \item How much of the bias is mechanical, the result of the tick, the fees and the cost of
  capital, and how much is preference, information or belief dispersion?
\end{enumerate}

We answer them with three data sets, each a stratified random sample with design weights
(\cref{sec:data}):
\begin{itemize}
  \item bid and ask quotes at fixed leads before resolution for 19,007 settled Kalshi markets in
  4,770 events, from 2021 to 2026;
  \item 2.69 million public Kalshi trade prints with size, price and aggressor side, covering
  15,329 markets in 18 strata;
  \item 2.57 million Polymarket fills by 237,569 identified wallets in 2,765 markets, and matched
  Kalshi and Polymarket prices for 441 games.
\end{itemize}

\paragraph{The bias is large per dollar, small per contract, and everywhere.} On Kalshi the
calibration slope of outcomes on prices is $1.13$ ($t=6.5$ against one). Longshots, sides priced
at 10\cents{} or less, return $-55\%$ per dollar at the mid-quote and $-75\%$ after the spread and
the taker fee. Favourites priced at 90\cents{} or more return $+1.1\%$ at the mid. The bias:
\begin{itemize}
  \item appears in every event category;
  \item does not fade as resolution approaches (the slope is 1.18 one hour ahead);
  \item is no weaker in the most liquid markets;
  \item has held in every quarter since 2025, while the venue grew by about three orders of
  magnitude.
\end{itemize}
It is steepest in multi-outcome events and ``by date'' deadline contracts. Pre-game sports prices
with a known start time are close to calibrated.

\paragraph{Small traders buy the longshots, and the price of immediacy falls on them.} In
Kalshi's trade prints:
\begin{itemize}
  \item 39\% of the contracts bought in taker orders of at most \$10 are longshots, against 1.2\%
  in orders above \$1,000. This holds in all 18 strata.
  \item In the median stratum, longshot positions lose for both liquidity roles. Takers lose
  more, in 13 of 18 strata,
  because they pay a half-spread and fee that together come to 14\% of a longshot's stake.
  \item Averaged over all prices, takers and makers earn about the same.
  \item The per-dollar longshot loss is similar at every order size, so small traders bear most
  of it because they hold most longshots, not because they trade at worse prices.
\end{itemize}
On Polymarket, where wallets are identified, retail wallets are the only significant net buyers of
longshots and account for 86\% of net longshot buying. No class of professional sellers takes the
other side: the flow is absorbed by many counterparties. Profits are extremely concentrated, with
the top 1\% of wallets taking 85\% of all gains, but wallet skill does not persist from one half of
the sample to the next.

\paragraph{Mechanics bound the bias; demand keeps prices at the bound.} Mechanics explain part of
the bias.
\begin{itemize}
  \item About half of the loss in the cheapest price bucket is the one-cent tick floor: contracts
  worth a tenth of a cent cannot be quoted lower.
  \item Averaged across markets, the longshot overpricing lies inside the no-arbitrage band set by
  the spread and the fee, so an arbitrageur taking the other side as a taker earns nothing.
\end{itemize}
But mechanics do not explain why prices sit near the edge of that band. The slope is unchanged
without the tick-floor buckets, and the bias does not shrink at short horizons, where capital
costs vanish. In the most heavily traded markets it exceeds the band.
An insider explanation fits poorly: Shin's insider share is small, and the field of longshots in
multi-outcome events stays overpriced after the Shin adjustment.
Proxies for belief dispersion, such as the number of outcomes, turnover and the presence of an
external forecast, do not predict where the bias is steeper. A single mild inverse-S
probability-weighting function, with the same curvature ($\gamma\approx0.89$) in every category,
describes prices well. Fitted on the first half of the sample, it removes 61\% of the longshot
overpricing in the second.
The evidence points to demand for lottery-like payoffs by small traders, met by liquidity
suppliers at prices that just cover their costs.

\paragraph{Contributions.} We contribute to three literatures.
\begin{itemize}
  \item \emph{The favourite--longshot bias.} We measure it on an order-book exchange across
  categories, horizons and four contract structures. We also separate the part that the tick and
  fees impose mechanically from the part that does not. Racetrack data rarely allow this
  \citep{snowbergwolfers2010,gandhi2015}.
  \item \emph{Prediction markets.} We extend \citet{burgi2025}, who document the bias in Kalshi
  trade prices and its stronger form for takers, to a volume-blind sample of all non-financial
  categories. We add complete trade histories classified by size and timing, and identified
  wallets on a second venue \citep{rothschildsethi2016,akey2026}.
  \item \emph{Retail lottery demand.} The longshot is the purest lottery ticket. The concentration
  of longshot buying in small accounts mirrors the evidence on lottery stocks and retail options
  trading \citep{kumar2009,bali2011,bryzgalova2023}. Here the claim has a fixed payoff date and no
  cash-flow risk, so its mispricing can be measured exactly.
\end{itemize}

The rest of the paper is organised as follows. \Cref{sec:lit} reviews related work and
\cref{sec:data} describes the venues and data. \Cref{sec:framework} sets out the measures, the
no-arbitrage band and the hypotheses. \Cref{sec:method} describes the methods and \cref{sec:results}
reports the results. \Cref{sec:discussion} discusses their interpretation and limits, and
\cref{sec:conclusion} concludes.
